% Section 4.4, from lectures/L04/appendix/c_exercises.md.
%
% Exercise 4 was titled "Över kanten", but with the course's layout (AB down, C across) its two
% groups, 0-2 and 5-7, are neighbours in the middle of the map and neither crosses an edge; it is
% "Från mintermer" here and now in the source too. Exercise 12 c), whose bold title runs into its
% sentence in the source, is a titled part here, with its sentence completed ("Bygg dem ...").

\section{Övningar}
\label{c4:sec:exercises}\label{c4:app:c}

\begin{aside}
\textbf{Så kontrollerar du ditt arbete.} Varje uttryck du tar fram med ett Karnaughdiagram kan du
kontrollera på två sätt, och du ska göra minst ett av dem varje gång: sätt in varje rad i
sanningstabellen i ditt uttryck och se att det ger rätt utgång, eller bygg nätet i CircuitVerse och
gå igenom alla kombinationer av ingångarna. Ett uttryck som inte är kontrollerat är en gissning.

Lösningsförslagen finns i bilaga~\ref{app:answers}. Gör varje övning innan du läser lösningen. Där
det finns ett vanligt felsvar visar lösningen det också, och varför det är fel.

Varje övning är märkt med sin sort. Kontrolluppgiften är övning~\ref{c4:ex:12}.
\end{aside}

\exercise{Graykod och grannar}{Förståelse}
\label{c4:ex:1}
\begin{parts}
  \item Varför skrivs rubrikerna i ett Karnaughdiagram i ordningen \code{00, 01, 11, 10} och inte
    \code{00, 01, 10, 11}? Svara med en mening om vad två grannrutor måste ha gemensamt.
  \item I ett diagram med fyra ingångar, \code{ABCD}: vilka fyra rutor är grannar till ruta 5
    (\code{ABCD = 0101})? Ange dem som mintermnummer. Använd figuren med rutornas nummer i
    \secref{c4:a2} (figur~\ref{c4:fig:layout4}) om du vill, men kontrollera svaret genom att se att
    varje granne skiljer sig från 0101 i exakt en bit.
  \item Samma fråga för ruta 0. Två av grannarna ligger på andra sidan en kant. Vilka?
  \item Varför får en grupp inte bestå av tre rutor, ens när tre ettor ligger på rad?
\end{parts}

\exercise{Två ingångar}{Räkna för hand}
\label{c4:ex:2}
\begin{narrowtable}{@{}ccc@{}}
  \thead{A} & \thead{B} & \thead{X} \\
  \midrule
  0 & 0 & 0 \\
  0 & 1 & 1 \\
  1 & 0 & 0 \\
  1 & 1 & 1 \\
\end{narrowtable}
\begin{parts}
  \item Rita Karnaughdiagrammet och ringa in grupperna.
  \item Läs av det minimerade uttrycket. Vad säger det om ingången \code{A}?
  \item Läs av uttrycket direkt ur tabellen, som i \secref{c2:a6}, och förenkla det algebraiskt. Får
    du samma svar?
\end{parts}

\exercise{Tre ingångar}{Räkna för hand}
\label{c4:ex:3}
\begin{narrowtable}{@{}cc@{}}
  \thead{ABC} & \thead{X} \\
  \midrule
  000 & 1 \\
  001 & 1 \\
  010 & 0 \\
  011 & 0 \\
  100 & 1 \\
  101 & 1 \\
  110 & 0 \\
  111 & 0 \\
\end{narrowtable}
\begin{parts}
  \item Fyll i diagrammet. Två av raderna med ettor ligger inte bredvid varandra på papperet. Är de
    grannar ändå?
  \item Ringa in och läs av det minimerade uttrycket.
  \item Hur många grindar behövs för det minimerade nätet, och hur många hade den direkta
    avläsningen ur tabellen krävt?
\end{parts}

\exercise{Från mintermer}{Räkna för hand}
\label{c4:ex:4}
\code{X} har tre ingångar \code{A}, \code{B}, \code{C} och är \code{1} för mintermerna 0, 2, 5
och 7.
\begin{parts}
  \item Skriv sanningstabellen, och fyll i diagrammet.
  \item Ringa in och läs av det minimerade uttrycket.
  \item Uttrycket beror bara på två av ingångarna. Vilken grind från kapitel~\ref{c2:ch} beräknar
    det med en enda grind?
\end{parts}

\exercise{Två utgångar}{Räkna för hand}
\label{c4:ex:5}
En krets har fyra ingångar \code{ABCD} och två utgångar \code{X} och \code{Y}:

\begin{narrowtable}{@{}cc@{\hspace{3em}}cc@{}}
  \thead{ABCD} & \thead{XY} & \thead{ABCD} & \thead{XY} \\
  \midrule
  0000 & 01 & 1000 & 11 \\
  0001 & 00 & 1001 & 10 \\
  0010 & 00 & 1010 & 10 \\
  0011 & 01 & 1011 & 11 \\
  0100 & 11 & 1100 & 01 \\
  0101 & 10 & 1101 & 00 \\
  0110 & 10 & 1110 & 00 \\
  0111 & 11 & 1111 & 01 \\
\end{narrowtable}
\begin{parts}
  \item Rita ett diagram för \code{X} och ett för \code{Y}, och läs av båda uttrycken.
  \item Två av ingångarna påverkar inte \code{X} alls, och två andra påverkar inte \code{Y}.
    Vilka? Hur syns det i diagrammen?
  \item Båda uttrycken kan skrivas med en enda grind var. Vilka grindar?
\end{parts}

\exercise{Hörnen}{Räkna för hand}
\label{c4:ex:6}
\code{X} har fyra ingångar och är \code{1} för mintermerna 0, 2, 8, 10 och 15.
\begin{parts}
  \item Fyll i diagrammet, ringa in och läs av.
  \item Minterm 15 får ingen granne att gå ihop med. Är det ett fel i din gruppering? Motivera.
\end{parts}

\exercise{Segment e i en sifferdisplay}{Räkna för hand}
\label{c4:ex:7}
En 7-segmentsdisplay visar en BCD-siffra 0--9 som kommer in på fyra ledningar \code{ABCD}
(\secref{c1:b5}). Segment \textbf{e}, det nedre vänstra, ska lysa för siffrorna 0, 2, 6 och 8, och
vara släckt för 1, 3, 4, 5, 7 och 9. Kombinationerna 10--15 förekommer aldrig.
\begin{parts}
  \item Fyll i diagrammet med ettor, och med X för 10--15.
  \item Ringa in och läs av ett minimerat uttryck för \code{e}. Använd X-rutorna där de gör
    grupperna större.
  \item Kontrollera ditt uttryck för alla tio siffrorna.
  \item Vad visar segment e om kretsen ändå får \code{1010} på ingången? Är det ett problem?
\end{parts}

\exercise{Fläkten}{Konstruktion}
\label{c4:ex:8}
En ventilationsfläkt \code{F} styrs av fyra signaler:
\begin{itemize}
  \item \code{T}: temperaturen är hög;
  \item \code{W}: fönstret är öppet;
  \item \code{C}: koldioxidhalten är hög;
  \item \code{B}: brandlarmet har löst ut.
\end{itemize}
Fläkten ska gå om temperaturen är hög och fönstret är stängt, eller om koldioxidhalten är hög. Men
när brandlarmet har löst ut ska fläkten \textbf{aldrig} gå, eftersom den då skulle mata elden med
luft.
\begin{parts}
  \item Skriv sanningstabellen, med ingångarna i ordningen \code{T W C B}.
  \item Fyll i Karnaughdiagrammet, med \code{TW} nedåt och \code{CB} åt sidan, och läs av ett
    minimerat uttryck.
  \item Rita grindnätet.
  \item Gör en grindtilldelning (\secref{c4:b7}): vilka kretsar ur 74HC08, 74HC32 och 74HC04
    behövs, och hur många av varje?
\end{parts}

\exercise{Enbart NAND}{Konstruktion}
\label{c4:ex:9}
\code{X = AB' + C}.
\begin{parts}
  \item Rita nätet med AND, OR och NOT.
  \item Gör om det till enbart NAND-grindar, enligt \secref{c4:a9}. Glöm inte inverteringen av
    \code{B}, och ingången \code{C} som går direkt till OR-grinden.
  \item Kontrollera med De Morgan att ditt NAND-uttryck är lika med \code{AB' + C}.
  \item Hur många NAND-grindar behövs, och hur många 74HC00?
\end{parts}

\exercise{Enbart NOR}{Konstruktion}
\label{c4:ex:10}
\code{X = (A + B)C}.
\begin{parts}
  \item Vilken form är uttrycket på, SP eller PS?
  \item Gör om det till enbart NOR-grindar. Rita nätet.
  \item Kontrollera ditt nät för alla åtta kombinationerna av \code{A}, \code{B} och \code{C}.
\end{parts}

\exercise{Kretsar och ben}{Förståelse}
\label{c4:ex:11}
Använd benplaceringarna i \secref{c4:b2}.
\begin{parts}
  \item Vilka ben har grind 3 i en 74HC00: ingångarna och utgången? Samma fråga för en 74HC02.
  \item Ett nät behöver sex AND-grindar och en inverterare. Hur många kretsar av vilka sorter?
  \item En grön lysdiod har spänningsfallet 2,1~V och ska drivas med ungefär 5~mA från en utgång
    på 5~V. Räkna ut förkopplingsmotståndet och välj ett standardvärde ur E12-serien (\dots, 330,
    390, 470, 560, 680, 820, 1000~Ω). Vilken ström blir det?
  \item En labbkamrat har lämnat ingångarna på de oanvända grindarna okopplade, och säger att det
    fungerar. Förklara varför det ändå är fel, och vad som kan hända.
\end{parts}

\exercise{Kontroll: minimerat och ominimerat nät}{Kontroll}
\label{c4:ex:12}
\emph{Räkna för hand, bygg och jämför, förklara.}

\code{X} har fyra ingångar och är \code{1} för mintermerna 1, 3, 5, 7, 9 och 11.

\task{a) För hand}
Skriv den direkta avläsningen ur tabellen: sex AND-termer med fyra ingångar var. Minimera sedan
med ett Karnaughdiagram.

\task{b) Förutsäg}
Hur många grindar och hur många kretsar behövs för vart och ett av de två uttrycken? Skriv ned
siffrorna innan du bygger något.

\task{c) Bygg båda i CircuitVerse}
Bygg dem med samma fyra ingångar, och varsin utgång. Gå igenom alla sexton kombinationerna och
anteckna båda utgångarna.

\task{d) Jämför}
Utgångarna ska vara lika på varje rad. Om de skiljer sig på någon rad: vilken rad, och vad säger
det om vilket av uttrycken som är fel?

\task{e)}
Det minimerade uttrycket kan skrivas \code{D(A' + B')}, och \code{A' + B'} är enligt De Morgan
\code{(AB)'}. Hur få grindar och kretsar klarar du dig med då? Jämför med ditt svar på b).

\exercise{Multiplexern}{Konstruktion, Fördjupning}
\label{c4:ex:13}
En \textbf{multiplexer} väljer vilken av två signaler som ska släppas fram. Den har två dataingångar
\code{A} och \code{B}, en väljare \code{S} och en utgång \code{X}: när \code{S = 0} ska
\code{X = B}, och när \code{S = 1} ska \code{X = A}.
\begin{parts}
  \item Skriv sanningstabellen med ingångarna i ordningen \code{S A B}.
  \item Minimera med ett Karnaughdiagram, med \code{SA} nedåt och \code{B} åt sidan.
  \item Rita nätet, och förklara med egna ord varför högst en av de två AND-termerna kan vara 1 åt
    gången.
\end{parts}
